\chapter{Power Series and Analytic Functions}
\label{ch:iv03}

Power series are the local algebra of holomorphic functions. Their radius of convergence controls a disk on which termwise differentiation and integration are valid, and analytic functions inherit remarkable rigidity.

\section*{Learning goals}
After this chapter, the reader should be able to:
\begin{itemize}
\item compute radii of convergence for power series and determine behavior at boundary points separately;
\item differentiate and integrate power series term by term inside their disks of convergence;
\item derive Taylor coefficients of holomorphic functions from derivatives at the center;
\item use local power-series expansions to prove analyticity and uniqueness statements;
\item manipulate standard expansions to obtain new series representations;
\item identify the nearest singularity controlling the radius of convergence in concrete examples;
\end{itemize}
\section*{Conceptual roadmap}
\[
\boxed{\text{local holomorphic structure}\;\longrightarrow\;\text{integral or series representation}\;\longrightarrow\;\text{global consequence}.}
\]

\section{Complex power series}
A series $\sum a_n(z-z_0)^n$ has a radius of convergence determined by coefficient growth.

% BEGIN VOL04-EXPANSION IV03-example-01
\begin{example}[Radius from coefficient growth]\label{ex:iv03-expand-01}
For
\[
\sum_{n=0}^{\infty}\frac{z^n}{3^n},
\]
the ratio of successive absolute terms tends to \(|z|/3\). Thus the series converges for \(|z|<3\) and diverges for \(|z|>3\), so its radius is \(R=3\). Inside the disk it sums to \(1/(1-z/3)\), but the radius is determined before any closed form is used.
\end{example}
% END VOL04-EXPANSION IV03-example-01
\section{Absolute convergence}
Inside the radius the series converges absolutely; outside it diverges.

\section{Uniform convergence on compact subdisks}
On every smaller closed disk, convergence is uniform and normally controlled.

\section{Termwise differentiation}
Differentiated power series has the same radius of convergence and represents the derivative.

% BEGIN VOL04-EXPANSION IV03-example-02
\begin{example}[Taylor expansion about a shifted center]\label{ex:iv03-expand-02}
To expand \(1/z\) about \(z_0=2\), write
\[
\frac1z=\frac1{2+(z-2)}=\frac12\frac1{1+(z-2)/2}
=\frac12\sum_{n=0}^{\infty}(-1)^n\left(\frac{z-2}{2}\right)^n.
\]
The nearest singularity to the center is at \(0\), at distance \(2\), so the expansion is valid exactly for \(|z-2|<2\).
\end{example}
% END VOL04-EXPANSION IV03-example-02
\section{Termwise integration}
Primitive series are obtained coefficientwise on disks.

\section{Analytic functions}
A function is analytic when locally represented by a convergent power series.

\section{Taylor coefficients}
For analytic functions, coefficients are derivatives divided by factorials.

% BEGIN VOL04-EXPANSION IV03-example-03
\begin{example}[Termwise differentiation keeps the radius]\label{ex:iv03-expand-03}
Starting from \(e^z=\sum_{n\ge0}z^n/n!\), termwise differentiation gives
\[
\sum_{n\ge1}\frac{nz^{n-1}}{n!}=\sum_{m\ge0}\frac{z^m}{m!}=e^z.
\]
The factorial growth makes the radius infinite both before and after differentiation. This concrete calculation models the general theorem that differentiation preserves the radius of a power series.
\end{example}
% END VOL04-EXPANSION IV03-example-03
\section{Uniqueness of expansion}
A power series representation on a disk is unique because its coefficients are recovered by derivatives at the center.

\section{Core structural results}

\begin{theorem}[Uniform convergence on compact subdisks]\label{thm:iv03-01}
If a power series has radius $R$, it converges uniformly on $|z-z_0|\le r<R$.
\begin{proof}
Choose $\rho$ with $r<\rho<R$. Convergence at radius $\rho$ bounds $|a_n|\rho^n$, giving a geometric majorant for $|a_n|r^n$.
\end{proof}
\end{theorem}

\begin{theorem}[Termwise differentiation]\label{thm:iv03-02}
Inside the radius of convergence, a power series may be differentiated termwise and the derivative series has the same radius.
\begin{proof}
Use uniform convergence on compact subdisks together with the standard derivative-series convergence criterion, or compare coefficients with the root test.
\end{proof}
\end{theorem}

\begin{theorem}[Coefficient uniqueness]\label{thm:iv03-03}
If $\sum a_n(z-z_0)^n=0$ on a neighborhood of $z_0$, then every $a_n=0$.
\begin{proof}
Differentiate $k$ times and evaluate at $z_0$ to get $k!a_k=0$.
\end{proof}
\end{theorem}

\section{Worked examples}

\begin{example}[Complex power series diagnostic]\label{ex:iv03-worked-01}
Give a rigorous worked analysis of Complex power series. State the relevant holomorphic, contour, series, or singularity hypothesis and derive the central conclusion. A series $\sum a_n(z-z_0)^n$ has a radius of convergence determined by coefficient growth. The decisive step is to use the complex-analytic structure globally rather than reason from real-variable intuition alone.
\end{example}

\begin{example}[Absolute convergence diagnostic]\label{ex:iv03-worked-02}
Give a rigorous worked analysis of Absolute convergence. State the relevant holomorphic, contour, series, or singularity hypothesis and derive the central conclusion. Inside the radius the series converges absolutely; outside it diverges. The decisive step is to use the complex-analytic structure globally rather than reason from real-variable intuition alone.
\end{example}

\begin{example}[Uniform convergence on compact subdisks diagnostic]\label{ex:iv03-worked-03}
Give a rigorous worked analysis of Uniform convergence on compact subdisks. State the relevant holomorphic, contour, series, or singularity hypothesis and derive the central conclusion. On every smaller closed disk, convergence is uniform and normally controlled. The decisive step is to use the complex-analytic structure globally rather than reason from real-variable intuition alone.
\end{example}

\begin{example}[Termwise differentiation diagnostic]\label{ex:iv03-worked-04}
Give a rigorous worked analysis of Termwise differentiation. State the relevant holomorphic, contour, series, or singularity hypothesis and derive the central conclusion. Differentiated power series has the same radius of convergence and represents the derivative. The decisive step is to use the complex-analytic structure globally rather than reason from real-variable intuition alone.
\end{example}

\section{Solved dossiers}
The dossiers are canonical solved problems. The provenance ledger separately records which retained corpus rules guided each topic and which dossiers were newly designed.

\begin{problem}[Complex power series diagnostic]\label{prob:iv03-dossier-01}
Give a rigorous worked analysis of Complex power series. State the relevant holomorphic, contour, series, or singularity hypothesis and derive the central conclusion.
\end{problem}
\begin{solution}
A series $\sum a_n(z-z_0)^n$ has a radius of convergence determined by coefficient growth. The decisive step is to use the complex-analytic structure globally rather than reason from real-variable intuition alone.
\end{solution}

\begin{problem}[Absolute convergence diagnostic]\label{prob:iv03-dossier-02}
Give a rigorous worked analysis of Absolute convergence. State the relevant holomorphic, contour, series, or singularity hypothesis and derive the central conclusion.
\end{problem}
\begin{solution}
Inside the radius the series converges absolutely; outside it diverges. The decisive step is to use the complex-analytic structure globally rather than reason from real-variable intuition alone.
\end{solution}

\begin{problem}[Uniform convergence on compact subdisks diagnostic]\label{prob:iv03-dossier-03}
Give a rigorous worked analysis of Uniform convergence on compact subdisks. State the relevant holomorphic, contour, series, or singularity hypothesis and derive the central conclusion.
\end{problem}
\begin{solution}
On every smaller closed disk, convergence is uniform and normally controlled. The decisive step is to use the complex-analytic structure globally rather than reason from real-variable intuition alone.
\end{solution}

\begin{problem}[Termwise differentiation diagnostic]\label{prob:iv03-dossier-04}
Give a rigorous worked analysis of Termwise differentiation. State the relevant holomorphic, contour, series, or singularity hypothesis and derive the central conclusion.
\end{problem}
\begin{solution}
Differentiated power series has the same radius of convergence and represents the derivative. The decisive step is to use the complex-analytic structure globally rather than reason from real-variable intuition alone.
\end{solution}

\begin{problem}[Termwise integration diagnostic]\label{prob:iv03-dossier-05}
Give a rigorous worked analysis of Termwise integration. State the relevant holomorphic, contour, series, or singularity hypothesis and derive the central conclusion.
\end{problem}
\begin{solution}
Primitive series are obtained coefficientwise on disks. The decisive step is to use the complex-analytic structure globally rather than reason from real-variable intuition alone.
\end{solution}

\begin{problem}[Analytic functions diagnostic]\label{prob:iv03-dossier-06}
Give a rigorous worked analysis of Analytic functions. State the relevant holomorphic, contour, series, or singularity hypothesis and derive the central conclusion.
\end{problem}
\begin{solution}
A function is analytic when locally represented by a convergent power series. The decisive step is to use the complex-analytic structure globally rather than reason from real-variable intuition alone.
\end{solution}

\begin{problem}[Taylor coefficients diagnostic]\label{prob:iv03-dossier-07}
Give a rigorous worked analysis of Taylor coefficients. State the relevant holomorphic, contour, series, or singularity hypothesis and derive the central conclusion.
\end{problem}
\begin{solution}
For analytic functions, coefficients are derivatives divided by factorials. The decisive step is to use the complex-analytic structure globally rather than reason from real-variable intuition alone.
\end{solution}

\begin{problem}[Uniqueness of expansion diagnostic]\label{prob:iv03-dossier-08}
Give a rigorous worked analysis of Uniqueness of expansion. State the relevant holomorphic, contour, series, or singularity hypothesis and derive the central conclusion.
\end{problem}
\begin{solution}
A power series representation on a disk is unique because its coefficients are recovered by derivatives at the center. The decisive step is to use the complex-analytic structure globally rather than reason from real-variable intuition alone.
\end{solution}

\begin{problem}[Uniform convergence on compact subdisks proof dossier]\label{prob:iv03-dossier-09}
Prove the following theorem-level statement and identify the essential hypothesis: If a power series has radius $R$, it converges uniformly on $|z-z_0|\le r<R$.
\end{problem}
\begin{solution}
Choose $\rho$ with $r<\rho<R$. Convergence at radius $\rho$ bounds $|a_n|\rho^n$, giving a geometric majorant for $|a_n|r^n$. This identifies the mechanism that makes the complex-analytic conclusion stronger than its real-variable analogue.
\end{solution}

\begin{problem}[Termwise differentiation proof dossier]\label{prob:iv03-dossier-10}
Prove the following theorem-level statement and identify the essential hypothesis: Inside the radius of convergence, a power series may be differentiated termwise and the derivative series has the same radius.
\end{problem}
\begin{solution}
Use uniform convergence on compact subdisks together with the standard derivative-series convergence criterion, or compare coefficients with the root test. This identifies the mechanism that makes the complex-analytic conclusion stronger than its real-variable analogue.
\end{solution}

\begin{problem}[Coefficient uniqueness proof dossier]\label{prob:iv03-dossier-11}
Prove the following theorem-level statement and identify the essential hypothesis: If $\sum a_n(z-z_0)^n=0$ on a neighborhood of $z_0$, then every $a_n=0$.
\end{problem}
\begin{solution}
Differentiate $k$ times and evaluate at $z_0$ to get $k!a_k=0$. This identifies the mechanism that makes the complex-analytic conclusion stronger than its real-variable analogue.
\end{solution}

\begin{problem}[Chapter synthesis]\label{prob:iv03-dossier-12}
Connect the definitions and principal theorems of this chapter into one reusable complex-analysis strategy.
\end{problem}
\begin{solution}
Power series are the local algebra of holomorphic functions. Their radius of convergence controls a disk on which termwise differentiation and integration are valid, and analytic functions inherit remarkable rigidity. The reusable strategy is to identify the analytic domain, choose the correct local expansion or contour representation, apply the structural theorem, and then audit singularities and topology.
\end{solution}

\section{Exercises with complete solutions}

\begin{exercise}\label{exr:iv03-01}
State and apply the central principle behind Complex power series. Give one concrete consequence.
\end{exercise}
\begin{hint}
Use the root or ratio test on the coefficients to find the radius; analyze points with \(|z-z_0|=R\) separately.
\end{hint}
\begin{solution}
A series $\sum a_n(z-z_0)^n$ has a radius of convergence determined by coefficient growth. The same principle can then be reused in later contour, residue, conformal, or Riemann-surface arguments.
\end{solution}

\begin{exercise}\label{exr:iv03-02}
State and apply the central principle behind Absolute convergence. Give one concrete consequence.
\end{exercise}
\begin{hint}
Inside a smaller closed disk, uniform convergence justifies termwise differentiation and integration.
\end{hint}
\begin{solution}
Inside the radius the series converges absolutely; outside it diverges. The same principle can then be reused in later contour, residue, conformal, or Riemann-surface arguments.
\end{solution}

\begin{exercise}\label{exr:iv03-03}
State and apply the central principle behind Uniform convergence on compact subdisks. Give one concrete consequence.
\end{exercise}
\begin{hint}
The Taylor coefficient is \(a_n=f^{(n)}(z_0)/n!\); compute derivatives only after locating the expansion center.
\end{hint}
\begin{solution}
On every smaller closed disk, convergence is uniform and normally controlled. The same principle can then be reused in later contour, residue, conformal, or Riemann-surface arguments.
\end{solution}

\begin{exercise}\label{exr:iv03-04}
State and apply the central principle behind Termwise differentiation. Give one concrete consequence.
\end{exercise}
\begin{hint}
Use the local series to identify the first nonzero coefficient; it controls local zero order and uniqueness.
\end{hint}
\begin{solution}
Differentiated power series has the same radius of convergence and represents the derivative. The same principle can then be reused in later contour, residue, conformal, or Riemann-surface arguments.
\end{solution}

\begin{exercise}\label{exr:iv03-05}
State and apply the central principle behind Termwise integration. Give one concrete consequence.
\end{exercise}
\begin{hint}
Rewrite the target into a known geometric, exponential, logarithmic, or binomial series before multiplying or composing expansions.
\end{hint}
\begin{solution}
Primitive series are obtained coefficientwise on disks. The same principle can then be reused in later contour, residue, conformal, or Riemann-surface arguments.
\end{solution}

\begin{exercise}\label{exr:iv03-06}
State and apply the central principle behind Analytic functions. Give one concrete consequence.
\end{exercise}
\begin{hint}
The nearest singularity in the complex plane bounds the Taylor disk even when the real-variable expression looks harmless.
\end{hint}
\begin{solution}
A function is analytic when locally represented by a convergent power series. The same principle can then be reused in later contour, residue, conformal, or Riemann-surface arguments.
\end{solution}

\begin{exercise}\label{exr:iv03-07}
State and apply the central principle behind Taylor coefficients. Give one concrete consequence.
\end{exercise}
\begin{hint}
To prove equality of two analytic expressions, compare their Taylor coefficients on a common disk of convergence.
\end{hint}
\begin{solution}
For analytic functions, coefficients are derivatives divided by factorials. The same principle can then be reused in later contour, residue, conformal, or Riemann-surface arguments.
\end{solution}

\begin{exercise}\label{exr:iv03-08}
State and apply the central principle behind Uniqueness of expansion. Give one concrete consequence.
\end{exercise}
\begin{hint}
Do not substitute outside the convergence disk; first determine the annular or disk region where each manipulated series is valid.
\end{hint}
\begin{solution}
A power series representation on a disk is unique because its coefficients are recovered by derivatives at the center. The same principle can then be reused in later contour, residue, conformal, or Riemann-surface arguments.
\end{solution}

\section*{Chapter summary}
The chapter now forms part of the canonical complex-analysis chain from holomorphicity through residues, global continuation, and Riemann surfaces.
% BEGIN VOL04-EXPANSION IV03-exercises-01
\section{Graded supplementary exercises}
These exercises extend the protected chapter with a balanced set of computations, proofs, hypothesis tests, applications, and challenges.

\subsection*{Standard computations}

\begin{exercise}[Geometric radius]\label{exr:iv03-expand-01}
Find the radius of convergence of \(\sum z^n/5^n\).
\end{exercise}
\begin{hint}
Apply the ratio test.
\end{hint}
\begin{solution}
The ratio tends to \(|z|/5\), so \(R=5\).
\end{solution}

\begin{exercise}[Factorial radius]\label{exr:iv03-expand-02}
Find the radius of convergence of \(\sum z^n/n!\).
\end{exercise}
\begin{hint}
Use the ratio test.
\end{hint}
\begin{solution}
The ratio of successive absolute terms is \(|z|/(n+1)\to0\) for every \(z\), hence \(R=\infty\).
\end{solution}

\begin{exercise}[Differentiate a series]\label{exr:iv03-expand-03}
Differentiate \(\sum_{n\ge0}z^n\) termwise inside its disk of convergence.
\end{exercise}
\begin{hint}
The geometric series has radius one.
\end{hint}
\begin{solution}
For \(|z|<1\), differentiation gives \(\sum_{n\ge1}nz^{n-1}=1/(1-z)^2\).
\end{solution}

\begin{exercise}[Taylor coefficients]\label{exr:iv03-expand-04}
Find the first four Taylor coefficients of \(e^z\) at zero.
\end{exercise}
\begin{hint}
Use \(a_n=f^{(n)}(0)/n!\).
\end{hint}
\begin{solution}
They are \(1,1,1/2,1/6\).
\end{solution}

\begin{exercise}[Shifted geometric series]\label{exr:iv03-expand-05}
Expand \(1/(3-z)\) in powers of \(z-1\) and state the radius.
\end{exercise}
\begin{hint}
Write \(3-z=2-(z-1)\).
\end{hint}
\begin{solution}
One obtains \(\frac12\sum_{n\ge0}((z-1)/2)^n\), valid for \(|z-1|<2\).
\end{solution}

\subsection*{Proofs}

\begin{exercise}[Absolute convergence inside the radius]\label{exr:iv03-expand-06}
Prove that a power series converges absolutely at every point strictly inside its radius.
\end{exercise}
\begin{hint}
Choose a larger radius where the coefficient terms are bounded.
\end{hint}
\begin{solution}
If \(|z-z_0|<r<R\), convergence at radius \(r\) bounds \(|a_n|r^n\). Multiplying by the geometric factor \((|z-z_0|/r)^n\) gives a summable majorant.
\end{solution}

\begin{exercise}[Uniform convergence on smaller disks]\label{exr:iv03-expand-07}
Prove uniform convergence on \(|z-z_0|\le r<R\).
\end{exercise}
\begin{hint}
Use the same geometric majorant independently of \(z\).
\end{hint}
\begin{solution}
Choose \(\rho\) with \(r<\rho<R\). Boundedness of \(|a_n|\rho^n\) yields the uniform majorant \(M(r/\rho)^n\), so the Weierstrass test applies.
\end{solution}

\begin{exercise}[Uniqueness of coefficients]\label{exr:iv03-expand-08}
Show that a power series representation near \(z_0\) has unique coefficients.
\end{exercise}
\begin{hint}
Differentiate repeatedly and evaluate at the center.
\end{hint}
\begin{solution}
Termwise differentiation gives \(f^{(n)}(z_0)=n!a_n\). Hence \(a_n=f^{(n)}(z_0)/n!\), fixing every coefficient.
\end{solution}

\begin{exercise}[Derivative radius]\label{exr:iv03-expand-09}
Explain why the derivative series has the same radius as the original series.
\end{exercise}
\begin{hint}
Compare \(|na_n|^{1/n}\) with \(|a_n|^{1/n}\).
\end{hint}
\begin{solution}
Since \(n^{1/n}\to1\), the limsup controlling the Cauchy--Hadamard radius is unchanged by multiplying coefficients by \(n\).
\end{solution}

\subsection*{Counterexamples and hypothesis tests}

\begin{exercise}[Boundary is separate]\label{exr:iv03-expand-10}
Does the radius of convergence determine convergence at every boundary point?
\end{exercise}
\begin{hint}
Compare the geometric series with an alternating harmonic type series.
\end{hint}
\begin{solution}
No. The radius separates the open convergence and exterior divergence regions, but individual points on the boundary require separate analysis.
\end{solution}

\begin{exercise}[Nearest singularity]\label{exr:iv03-expand-11}
Why can the Taylor series of \(1/(1-z)\) at zero not have radius larger than one?
\end{exercise}
\begin{hint}
The represented analytic function has a singularity at \(z=1\).
\end{hint}
\begin{solution}
A convergent power series is analytic throughout its disk. A radius larger than one would make the series analytic at \(1\), contradicting the pole there.
\end{solution}

\begin{exercise}[Smooth is not analytic]\label{exr:iv03-expand-12}
Does infinite real differentiability imply a complex power series representation?
\end{exercise}
\begin{hint}
Complex analyticity is much more rigid than real smoothness.
\end{hint}
\begin{solution}
No. Real smoothness alone imposes no Cauchy--Riemann structure. A function of \(x\) and \(y\) can be smooth as a real map but fail complex differentiability everywhere.
\end{solution}

\subsection*{Applications and investigations}

\begin{exercise}[Summing a weighted geometric series]\label{exr:iv03-expand-13}
Use power series to evaluate \(\sum_{n\ge1}nr^{n-1}\) for \(|r|<1\).
\end{exercise}
\begin{hint}
Differentiate the geometric identity.
\end{hint}
\begin{solution}
Differentiating \(\sum r^n=1/(1-r)\) gives \(\sum_{n\ge1}nr^{n-1}=1/(1-r)^2\).
\end{solution}

\begin{exercise}[Local analytic model]\label{exr:iv03-expand-14}
Explain how a Taylor polynomial approximates an analytic function near its center.
\end{exercise}
\begin{hint}
Separate the first \(N\) terms from the convergent tail.
\end{hint}
\begin{solution}
Inside a smaller disk the tail is uniformly controlled, so truncating after degree \(N\) gives a computable local polynomial model whose error tends to zero as \(N\) grows.
\end{solution}

\subsection*{Challenge problems}

\begin{exercise}[Cauchy--Hadamard]\label{exr:iv03-expand-15}
State the coefficient formula for the radius of convergence and explain its meaning.
\end{exercise}
\begin{hint}
Use the root test on \(|a_n(z-z_0)^n|\).
\end{hint}
\begin{solution}
The formula is \(1/R=\limsup |a_n|^{1/n}\), with the usual conventions. It measures the exponential growth rate of the coefficients.
\end{solution}

\begin{exercise}[Multiple centers]\label{exr:iv03-expand-16}
The function \(1/(1-z)\) is expanded about \(z_0=-1\). Find the radius without computing coefficients.
\end{exercise}
\begin{hint}
Measure the distance from the center to the nearest singularity.
\end{hint}
\begin{solution}
The only finite singularity is at \(1\), distance \(2\) from \(-1\). Therefore the Taylor radius is \(2\).
\end{solution}

% END VOL04-EXPANSION IV03-exercises-01
